\documentclass[11pt]{article}
\usepackage[margin=1in]{geometry}
\usepackage{amsmath, amssymb, amsthm}
\usepackage{booktabs}
\usepackage{enumitem}
\usepackage{graphicx}
\usepackage{url}

\newtheorem{theorem}{Theorem}
\newtheorem{lemma}[theorem]{Lemma}
\newtheorem{proposition}[theorem]{Proposition}
\newtheorem{remark}[theorem]{Remark}
\theoremstyle{definition}
\newtheorem{definition}[theorem]{Definition}

\title{The Krenn--Gu conjecture holds for all simple graphs\\[0.3em]
on six vertices}
\author{Junhao Liang\\[0.4em]
Zhejiang Normal University\\[0.2em]
\url{mailto:JunhaoLiang2005@gmail.com}}
\date{}

\begin{document}
\maketitle

\begin{abstract}
We prove that for \emph{every} number of colors $D\ge 3$ there is no
bi-colored weighted \emph{simple} graph on $6$ vertices --- bichromatic
edges and complex weights allowed --- in which all $D$ monochromatic
inherited vertex colorings have unit weight while every other vertex
coloring has weight zero. Equivalently, the matching index of every
simple graph on $6$ vertices is at most $2$; $K_4$ (with $\mu=3$)
remains the only known exception to $\mu(G)\le 2$ at this order. For
$D=3,4,5$ the proof reduces the statement to a purely combinatorial
finite search: a complete enumeration of the $5{,}610$, $2{,}160$ and
$720$ possible monochromatic-class configurations respectively,
collapsed to $24$, $5$ and $1$ orbit representatives under the symmetric
group $S_6$, followed by an exhaustive check of all $2{,}033{,}610$,
$4{,}917$ and $1$ edge-color assignments --- every one of which is
eliminated by one of two elementary necessary conditions; for $D\ge 6$
the statement follows from a one-line edge-counting argument
($3D\le 15$). Because the elimination is weight-free, the result holds
over any integral domain, in particular over $\mathbb C$, $\mathbb R$,
$\mathbb Z$, and for weights in $\{-1,0,1\}$. The certificate is
machine-reproducible (seconds of computation), has been independently
re-verified by fresh implementations and by an independent SAT encoding,
and settles the open statements
\texttt{eqSystem6\_no\_solution\_d3}, \texttt{\_d4}, \texttt{\_d5} and
\texttt{\_ge3} (with their \texttt{\_real}, \texttt{\_int},
\texttt{\_trinary\_int} variants where present) in the DeepMind
\emph{formal-conjectures} benchmark. The full conjecture (arbitrary
$n\ge 8$, multigraphs) remains open.
\end{abstract}

\section{Introduction}\label{sec:intro}

The following conjecture emerged from quantum optics in 2017
\cite{KGZ17} and was stated in graph-theoretic form by Krenn, Gu and
Solt\'esz \cite{KGS19}:

\begin{quote}
\emph{If $G$ is a graph with more than $4$ vertices, then the matching
index satisfies $\mu(G)\le 2$.}
\end{quote}

Here $\mu(G)$ is the dimension of the subspace spanned by
``monochromatic'' perfect-matchings superpositions; the complete graph
$K_4$ attains $\mu(K_4)=3$ and is conjecturally the only graph of
matching index at least $3$ (see \cite{KGS19, GKM24} for the exact
definitions and the prize page \cite{bounty} for the current status and
the prize associated with a full resolution). The conjecture has
survived substantial computational effort: SAT-based searches ruled out
mono-edge constructions for small orders \cite{CLA22}, and the
best partial results show $\mu(G)\le 2$ for graphs of vertex
connectivity at most $2$ and for cubic graphs, with any minimal
counterexample forced to be $4$-connected \cite{CGI24}. For
\emph{simple} graphs, the strongest general bound is
$\mu(G)\le n/\sqrt{2}$ (Theorem~22 of \cite{CG24}); at $n=6$ this yields
only $\mu(G)\le 4$, since $\mu$ is integral, whereas the present paper
proves $\mu(G)\le 2$. Over
\emph{nonnegative} real weights the conjecture is fully settled
(Bogdanov, see \cite{mathoverflow}), and the DeepMind formal-conjectures
benchmark records the general case as open, with only the $D=N$ cases
resolved formally \cite{alphaproof, formalconjectures}.

In this note we resolve the case $n=6$ \emph{completely for simple
graphs}: the conjecture holds for every number of colors $d\ge 3$.
Our main result is:

\begin{theorem}\label{thm:main}
For every $D\ge 3$ there is no bi-colored weighted simple graph on $6$
vertices --- with each edge carrying an ordered color pair in
$\{0,\dots,D-1\}^2$ and a nonzero weight in an integral domain $R$ ---
such that $w(\sigma)=1$ for every monochromatic coloring
$\sigma=i^6$, $i\in\{0,\dots,D-1\}$, and $w(\sigma)=0$ for every other
vertex coloring.
\end{theorem}

Consequently $\mu(G)\le 2$ for every simple graph $G$ on $6$ vertices
(over any integral domain). Two features deserve emphasis:

\begin{enumerate}[label=(\alph*)]
\item The obstruction is \emph{purely combinatorial}: the necessary
conditions (b) and (c) of Proposition~\ref{prop:necessary} below never
mention weights. Hence the result is independent of the coefficient
domain.
\item The certificate is \emph{exhaustive and machine-checkable}: the
entire search space ($2{,}033{,}610$ color assignments) is enumerated
and eliminated in seconds, with no reliance on symbolic or numeric
solvers.
\end{enumerate}

The result settles the DeepMind formal-conjectures open statements
\texttt{eqSystem6\_no\_solution\_d3}, \texttt{\_d4}, \texttt{\_d5} and
the universal \texttt{\_ge3} family, with their
\texttt{\_real}, \texttt{\_int} and \texttt{\_trinary\_int} variants
where present (Section~\ref{sec:consequences}).

\section{Preliminaries}\label{sec:prelim}

We follow the terminology of \cite{KGS19, CGI24}. Fix the vertex set
$V=\{0,1,2,3,4,5\}$. A \emph{bi-colored weighted graph} is a triple
$(V,E,c,w)$ where $E\subseteq\binom{V}{2}$ is a simple graph, each edge
$e=\{u,v\}\in E$ ($u<v$) carries an ordered color pair
$c(e)=(c_1,c_2)\in\{0,1,2\}^2$ ($c_1$ at $u$, $c_2$ at $v$), and
$w:E\to R$ assigns to each edge a nonzero weight in a commutative
integral domain $R$. A \emph{perfect matching} (PM) of $(V,E)$ is a set
of three pairwise disjoint edges covering all six vertices. For a PM
$M$, the \emph{inherited vertex coloring} $\mathrm{IVC}(M)\in
\{0,1,2\}^6$ is defined by giving each vertex the color of its incident
matching edge at that vertex; for $\sigma\in\{0,1,2\}^6$ put
\[
w(\sigma)\;=\;\sum_{M:\ \mathrm{IVC}(M)=\sigma}\ \prod_{e\in M} w(e).
\]
A \emph{$(6,3)$-solution} is a bi-colored weighted graph with
$w(000000)=w(111111)=w(222222)=1$ and $w(\sigma)=0$ for all other
$\sigma$.

The fifteen PMs of $K_6$ are the candidates; a PM $M$ of $K_6$
\emph{survives} in $(V,E,c,w)$ if $M\subseteq E$.

\section{Reduction to a finite combinatorial search}\label{sec:reduce}

The following lemmas are elementary; they distill the argument of
\cite{KGS19} and the project notes \cite{notes}. All of them hold over
any integral domain.

\begin{lemma}[All-nonzero WLOG]\label{lem:nonzero}
If a $(6,3)$-solution exists, one exists with all weights nonzero.
\end{lemma}
\begin{proof}
Delete every edge of weight $0$: such edges contribute $0$ to every
surviving PM product, so every sum $w(\sigma)$ is unchanged, and the
monochromatic sums remain $1$, so each color still has a surviving PM.
\end{proof}

\begin{lemma}[$K_6$ suffices]\label{lem:k6}
If some simple graph on $6$ vertices has a $(6,3)$-solution, then so
does $K_6$ (with missing edges colored arbitrarily and given weight
$0$).
\end{lemma}
\begin{proof}
Extend the coloring of $E$ to all fifteen pairs arbitrarily and set the
weight of every new edge to $0$. Every new PM uses a zero-weight edge
and contributes nothing.
\end{proof}

By Lemmas~\ref{lem:nonzero} and \ref{lem:k6}, the search space is:
color each of the fifteen pairs of $V$ by an ordered pair in
$\{0,1,2\}^2$ or mark it \emph{absent} (weight $0$), and keep only
colorings in which every surviving PM has all nonzero weights --- which,
by Lemma~\ref{lem:nonzero}, every surviving solution may be assumed to
satisfy.

\begin{lemma}[Monochromatic PMs]\label{lem:mono}
For a PM $M$ with all edges present,
$\mathrm{IVC}(M)=iiiiii$ (write $i^6$) iff every edge of $M$ has color
pair $(i,i)$.
\end{lemma}
\begin{proof}
Vertex $u$ receives color $i$ iff its matching edge has color $i$ at
$u$. If all six vertices receive color $i$, both endpoint colors of
every matching edge equal $i$.
\end{proof}

Write $P_i$ for the set of surviving PMs with IVC $i^6$; by
Lemma~\ref{lem:mono}, $P_i=\mathrm{PM}(H_i)$ where $H_i$ is the subgraph
of color-$(i,i)$ edges. An edge has a single color pair, so $H_0,H_1,H_2$
are pairwise edge-disjoint.

\begin{lemma}[Per-IVC decomposition]\label{lem:decomp}
$w(\sigma)$ depends only on the sets $P_i$ and the non-monochromatic
IVC groups: $w(i^6)=\sum_{M\in P_i}\prod_{e\in M}w(e)$, and
$w(\sigma)=\sum_{M:\,\mathrm{IVC}(M)=\sigma}\prod_{e\in M}w(e)$ for
non-monochromatic $\sigma$.
\end{lemma}

\begin{lemma}[The L3 condition]\label{lem:L3}
In a solution with all weights nonzero, every \emph{realized}
non-monochromatic IVC is realized by at least two surviving PMs.
\end{lemma}
\begin{proof}
If a unique surviving PM $M$ realizes $\sigma$, then
$w(\sigma)=\prod_{e\in M}w(e)\neq 0$ (nonzero product in an integral
domain), contradicting $w(\sigma)=0$.
\end{proof}

\begin{proposition}[Necessary conditions]\label{prop:necessary}
Every $(6,3)$-solution induces a coloring of the fifteen pairs of
$V$ --- each pair colored by a pair in $\{0,1,2\}^2$ or absent ---
satisfying
\begin{enumerate}[label=(\alph*)]
\item $P_0,P_1,P_2$ are all nonempty (Lemma~\ref{lem:decomp} with
$w(i^6)=1$),
\item no surviving PM outside $P_0\cup P_1\cup P_2$ has a monochromatic
IVC (Lemma~\ref{lem:mono}),
\item every realized non-monochromatic IVC is realized by at least two
surviving PMs (Lemma~\ref{lem:L3}).
\end{enumerate}
\end{proposition}

It is (a)--(c) --- and nothing else --- that the search below
eliminates. Note that no condition on the coefficient domain $R$ was
used anywhere; the whole argument is weight-free.

\section{The exhaustive certificate}\label{sec:search}

\subsection{Class enumeration and orbit reduction}

By (a) and Lemma~\ref{lem:mono}, a surviving solution determines an
ordered triple $(P_0,P_1,P_2)$ of \emph{nonempty}, \emph{pairwise
edge-disjoint} subsets of the set of the fifteen PMs of $K_6$. A direct
construction enumerates all such triples:

\begin{proposition}[Class count, computational]\label{prop:5610}
There are exactly $5{,}610$ ordered triples $(P_0,P_1,P_2)$ of
nonempty, pairwise edge-disjoint PM subsets of $K_6$.
\end{proposition}

The symmetric group $S_6$ acts on vertex labels; we fix the three color
labels throughout (no color renaming is used). Canonicalizing each
triple under this action collapses the $5{,}610$ classes to $24$
orbit representatives, listed in Appendix~\ref{app:reps}. The
completeness statement is:

\begin{proposition}[Orbit completeness, computational]\label{prop:24}
For every one of the $5{,}610$ triples of
Proposition~\ref{prop:5610}, some permutation in $S_6$ sends it to one
of the $24$ representatives in Appendix~\ref{app:reps}; the $24$
representatives are pairwise inequivalent.
\end{proposition}

\subsection{The batch}

Fix an orbit representative $(P_0,P_1,P_2)$. The edges appearing in the
PMs of $P_0\cup P_1\cup P_2$ are \emph{forced}: an edge of a PM of
$P_i$ must carry the color pair $(i,i)$ (Lemma~\ref{lem:mono}). Every
remaining (\emph{free}) edge may be colored by any of the $9$ ordered
pairs in $\{0,1,2\}^2$ or be absent. Writing $F$ for the number of free
edges of the representative, the batch enumerates all $10^{F}$
assignments. For every assignment we compute the IVC code of each of
the fifteen PMs of $K_6$ (a PM using an absent edge is \emph{dead} and
contributes nothing) and test the two conditions:

\begin{description}
\item[no-new-mono (b):] no PM outside its class (a dead PM, or a PM of
$P_j$ for $j\neq i$, or an unclassified PM) has IVC $i^6$ for any
$i\in\{0,1,2\}$;
\item[L3 (c):] among the \emph{realized} non-monochromatic IVC codes,
no code occurs exactly once.
\end{description}

A coloring that passes both tests would proceed to the exact decision
procedure (linear system over $\mathbb Q$ together with the
liftability relations $\ker U^{\mathsf T}$); see \cite{notes}. This
stage is never reached:

\begin{theorem}[Exhaustive elimination, computational]\label{thm:search}
For each of the $24$ representatives of Appendix~\ref{app:reps},
every one of the $10^{F}$ assignments fails (b) or (c). In total,
$2{,}033{,}610$ assignments are eliminated; the number of survivors
is $0$.
\end{theorem}

Table~\ref{tab:classes} records the distribution of the
representatives by class sizes and free-edge count; the total is
$\sum 10^{F}=2{,}033{,}610$.

\begin{table}[h]
\centering
\begin{tabular}{lll}
\toprule
class size $(|P_0|,|P_1|,|P_2|)$ & free edges $F$ & multiplicity \\
\midrule
$(1,1,1)$ & $6$ & $2$ \\
$(1,1,2)$ & $3,4$ & $1,1$ \\
$(1,1,3)$ & $0,2$ & $1,1$ \\
$(1,1,4)$ & $0$ & $1$ \\
$(1,2,1)$ & $4,3$ & $1,1$ \\
$(1,2,2)$ & $2,0$ & $1,1$ \\
$(1,3,1)$ & $2,0$ & $1,1$ \\
$(1,4,1)$ & $0$ & $1$ \\
$(2,1,1)$ & $4,3$ & $1,1$ \\
$(2,1,2)$ & $2,0$ & $1,1$ \\
$(2,2,1)$ & $2,0$ & $1,1$ \\
$(2,2,2)$ & $0$ & $1$ \\
$(3,1,1)$ & $2,0$ & $1,1$ \\
$(4,1,1)$ & $0$ & $1$ \\
\bottomrule
\end{tabular}
\caption{The $24$ orbit representatives: class sizes, free-edge counts,
and multiplicities. Total rows: $2\cdot 10^{6}+10^{4}\cdot 2+
10^{3}\cdot 2+10^{2}\cdot 4+1\cdot 8=2{,}033{,}610$.}
\label{tab:classes}
\end{table}

\emph{Proof of Theorem~\ref{thm:main}.} Suppose a $(6,3)$-solution
exists over an integral domain $R$. By
Lemmas~\ref{lem:nonzero}--\ref{lem:k6} it induces a coloring of the
fifteen pairs (each pair colored or absent) with all necessary weights
nonzero, whose triple $(P_0,P_1,P_2)$ is one of the $5{,}610$ classes.
By Proposition~\ref{prop:24} some $S_6$-image of this coloring has the
same triple as one of the $24$ representatives, and the induced free-
edge assignment is one of the $10^{F}$ enumerated. By
Proposition~\ref{prop:necessary} the coloring satisfies (a)--(c); but
Theorem~\ref{thm:search} says no enumerated assignment satisfies (b)
and (c) --- contradiction. \qed

\subsection{Four and five colors, and the case $D\ge 6$}

The method extends verbatim to $D=4,5$ colors. A solution with $D$
colors determines $D$ pairwise edge-disjoint nonempty monochromatic PM
classes, and the edge-counting argument below already bounds the
possibilities:

\begin{lemma}[Only five colors can occur]\label{lem:Dle5}
In a $(6,D)$-solution over any integral domain, $D\le 5$.
\end{lemma}
\begin{proof}
For each color $i$ the monochromatic sum $w(i^6)=1$ forces the class
$P_i$ to be nonempty (an empty sum is $0\neq 1$). The classes are
pairwise edge-disjoint (an edge carries one color pair), and each
nonempty class contains a PM, hence at least $3$ edges. Therefore
$3D\le 15$.
\end{proof}

For $D=4$: enumerating all ordered quadruples of nonempty pairwise
edge-disjoint PM subsets yields $2{,}160$ classes, which collapse to
$5$ orbit representatives under $S_6$ (Appendix~\ref{app:reps}); every
free edge has $16+1$ options, the free-edge count satisfies
$F\le 15-12=3$, and the complete batch enumerates
$4{,}917$ assignments, all eliminated by (b) or (c).

For $D=5$: an ordered quintuple forces all five classes to be single
PMs (two PMs in one class would force the union to contain
$5+12=17>15$ edges), i.e.\ the classes form an ordered
$1$-factorization of $K_6$; there are exactly $6\cdot 5!=720$ such
quintuples, all equivalent under $S_6$ --- one orbit, one forced
coloring ($F=0$), eliminated by (b) or (c).

Combined with Lemma~\ref{lem:Dle5}, the three exhaustive searches settle
Theorem~\ref{thm:main} in full.

\section{Independent verification}\label{sec:verify}

Because the theorem rests on a computation, the computation was
re-implemented from scratch and cross-validated by independent means.
All artifacts (scripts, logs, this certificate) accompany the paper
\cite{repo}.

\begin{enumerate}[label=(V\arabic*)]
\item \emph{Fresh implementation.} The PM enumeration, the triple
enumeration, the canonicalization, the IVC computation, and the filter
were rewritten independently (no shared decision logic), reproducing
$5{,}610$ and the $24$-orbit completeness claim, and re-running the
full batch: $2{,}033{,}610$ rows, $0$ survivors.
\item \emph{Filter validation.} The vectorized filter was cross-checked
against a naive per-row reference implementation on $20{,}000$ random
rows, and the original implementation's filter was cross-checked
against the same reference on $60{,}000$ random rows; full agreement.
\item \emph{SAT confirmation.} An independent encoding of the
constraints (b) and (c) for each of the $24$ classes was solved with
the SMT solver z3: $\textsf{UNSAT}$ in all $24$ cases.
\item \emph{Decision-procedure controls.} A second, independently
written exact decision engine agrees with the original engine on
$300$ random colorings and $200$ structured instances ($100\%$
agreement), and passes positive controls: the known $3$-dimensional
example on $4$ vertices (the $K_4$ witness) is correctly certified
feasible with an exact weight construction, and relaxed instances
exercise the lift-relation machinery.
\item \emph{Mono-edge subcase.} Of the $3^{15}=14{,}348{,}907$
mono-edge colorings, exactly $248{,}310$ realize all three
monochromatic IVCs, and $0$ of them satisfy (c) --- recovering and
strengthening the SAT result of \cite{CLA22} for $n=6$.
\item \emph{Four and five colors.} The $D=4$ and $D=5$ enumerations
($2{,}160$ and $720$ classes; $5$ and $1$ orbits; $4{,}917$ and $1$
rows; $0$ survivors) were re-derived from scratch: the raw counts by an
independent formula ($\sum_T(2^{c(T)}-1)$ over the audited triple list,
resp.\ $6\cdot 5!$), the orbit cover by re-canonicalizing every raw
class, the batches replayed with an independent naive filter
implementation, and the constraints confirmed by z3 ($\textsf{UNSAT}$
$5/5$ and $1/1$).
\end{enumerate}

No randomness is involved in any infeasibility verdict: every
elimination in Theorem~\ref{thm:search} is a structural check of
condition (b) or (c).

\section{Consequences for the formal-conjectures benchmark}
\label{sec:consequences}

The DeepMind \emph{formal-conjectures} benchmark
\cite{formalconjectures, alphaproof} contains the open Lean statements
$\neg\,\exists\,W:\mathrm{WeightsN}\,6\,D\ R,\
\mathrm{EqSystemN}\,6\,D\ W$ for various color counts $D$ and
coefficient domains $R$. Since Theorem~\ref{thm:main} holds over any
integral domain, it settles the following $13$ statements currently
recorded as open in the benchmark:
\[
\begin{gathered}
\texttt{eqSystem6\_no\_solution\_d3},\quad
\texttt{\_d3\_real},\quad
\texttt{\_d3\_int},\quad
\texttt{\_d3\_trinary\_int},\\
\texttt{eqSystem6\_no\_solution\_d4},\quad
\texttt{eqSystem6\_no\_solution\_d5},\quad
\texttt{\_d5\_real},\quad
\texttt{\_d5\_int},\quad
\texttt{\_d5\_trinary\_int},\\
\texttt{eqSystem6\_no\_solution\_ge3},\quad
\texttt{\_ge3\_real},\quad
\texttt{\_ge3\_int},\quad
\texttt{\_ge3\_trinary\_int}.
\end{gathered}
\]
(The benchmark also records \texttt{eqSystem6\_no\_solution\_d6} as
solved; it follows elementarily from Lemma~\ref{lem:Dle5} as well.)

\section{Conclusion and open problems}\label{sec:open}

We have settled the Krenn--Gu conjecture for all simple graphs on $6$
vertices and all color counts $D\ge 3$. The following directions
remain:

\begin{enumerate}[label=(\roman*)]
\item \emph{Multigraphs.} The conjecture allows parallel edges. The
known reductions merge same-pair parallels and delete self-cancelling
families; whether \emph{distinct-pair} parallel edges can be
eliminated (thereby extending Theorem~\ref{thm:main} to multigraphs on
$6$ vertices) is open. Computational evidence (restricted layers,
cascade analysis) is recorded in the accompanying repository.
\item \emph{$n\ge 8$.} The combinatorial bottleneck (condition (c)) is
exactly what fails to generalize naively; new ideas are needed.
\item \emph{Formalization.} A Lean formalization of the three
enumerations (the $24$-, $5$- and $1$-orbit searches and their
$2{,}033{,}610$-, $4{,}917$- and $1$-row eliminations, executable under
the Lean kernel) would turn the $13$ benchmark statements of
Section~\ref{sec:consequences} into \texttt{answer(True)} entries.
\end{enumerate}

\begin{remark}[Relation to the prize]
The prize of \texteuro 3{,}000 \cite{bounty} is attached to a full
resolution of the conjecture (all even $n\ge 6$, multigraphs
included). This note does not claim it.
\end{remark}

\section*{Acknowledgments}
This research was conducted with substantial assistance from AI systems
(Anthropic's Claude), which contributed to the design of the exhaustive
search, the drafting of the manuscript, and the independent verification
described in Section~\ref{sec:verify}. All mathematical claims were
reviewed by the author, and every computational claim was re-derived by
at least one independent implementation or independent solver. No AI
system is a co-author of this paper.

\appendix

\section{The $24$ orbit representatives}\label{app:reps}

Vertices are $0,\dots,5$; each PM is written as three edges, e.g.\
$01\text{-}23\text{-}45$ denotes $\{\{0,1\},\{2,3\},\{4,5\}\}$.
Representative number $k$ has classes $(P_0^k, P_1^k, P_2^k)$.

\begin{center}
\resizebox{\textwidth}{!}{%
\begin{tabular}{clll}
\toprule
$k$ & $P_0^k$ & $P_1^k$ & $P_2^k$ \\
\midrule
$0$ & $01\text{-}23\text{-}45$ & $02\text{-}14\text{-}35$ & $03\text{-}15\text{-}24$ \\
$1$ & $01\text{-}23\text{-}45$ & $02\text{-}14\text{-}35$ & $03\text{-}15\text{-}24\ /\ 04\text{-}13\text{-}25$ \\
$2$ & $01\text{-}23\text{-}45$ & $02\text{-}14\text{-}35$ & $03\text{-}15\text{-}24\ /\ 04\text{-}13\text{-}25\ /\ 05\text{-}12\text{-}34$ \\
$3$ & $01\text{-}23\text{-}45$ & $02\text{-}14\text{-}35$ & $05\text{-}13\text{-}24$ \\
$4$ & $01\text{-}23\text{-}45$ & $02\text{-}14\text{-}35$ & $03\text{-}15\text{-}24\ /\ 05\text{-}13\text{-}24$ \\
$5$ & $01\text{-}23\text{-}45$ & $02\text{-}14\text{-}35$ & $03\text{-}15\text{-}24\ /\ 04\text{-}13\text{-}25\ /\ 05\text{-}13\text{-}24$ \\
$6$ & $01\text{-}23\text{-}45$ & $02\text{-}14\text{-}35$ & $03\text{-}15\text{-}24\ /\ 04\text{-}13\text{-}25\ /\ 05\text{-}12\text{-}34\ /\ 05\text{-}13\text{-}24$ \\
$7$ & $01\text{-}23\text{-}45$ & $02\text{-}14\text{-}35\ /\ 02\text{-}15\text{-}34$ & $04\text{-}13\text{-}25$ \\
$8$ & $01\text{-}23\text{-}45$ & $02\text{-}14\text{-}35\ /\ 02\text{-}15\text{-}34$ & $04\text{-}13\text{-}25\ /\ 05\text{-}13\text{-}24$ \\
$9$ & $01\text{-}23\text{-}45$ & $02\text{-}15\text{-}34\ /\ 03\text{-}14\text{-}25$ & $04\text{-}12\text{-}35$ \\
$10$ & $01\text{-}23\text{-}45$ & $02\text{-}15\text{-}34\ /\ 03\text{-}14\text{-}25$ & $04\text{-}12\text{-}35\ /\ 05\text{-}13\text{-}24$ \\
$11$ & $01\text{-}23\text{-}45$ & $02\text{-}14\text{-}35\ /\ 02\text{-}15\text{-}34\ /\ 03\text{-}14\text{-}25$ & $05\text{-}13\text{-}24$ \\
$12$ & $01\text{-}23\text{-}45$ & $02\text{-}15\text{-}34\ /\ 03\text{-}14\text{-}25\ /\ 04\text{-}12\text{-}35$ & $05\text{-}13\text{-}24$ \\
$13$ & $01\text{-}23\text{-}45$ & $02\text{-}14\text{-}35\ /\ 02\text{-}15\text{-}34\ /\ 03\text{-}14\text{-}25\ /\ 04\text{-}12\text{-}35$ & $05\text{-}13\text{-}24$ \\
$14$ & $01\text{-}23\text{-}45\ /\ 01\text{-}24\text{-}35$ & $02\text{-}15\text{-}34$ & $03\text{-}14\text{-}25$ \\
$15$ & $01\text{-}23\text{-}45\ /\ 01\text{-}24\text{-}35$ & $02\text{-}15\text{-}34$ & $03\text{-}14\text{-}25\ /\ 04\text{-}13\text{-}25$ \\
$16$ & $01\text{-}23\text{-}45\ /\ 01\text{-}24\text{-}35$ & $03\text{-}14\text{-}25\ /\ 04\text{-}13\text{-}25$ & $02\text{-}15\text{-}34$ \\
$17$ & $01\text{-}23\text{-}45\ /\ 01\text{-}24\text{-}35$ & $03\text{-}14\text{-}25\ /\ 04\text{-}13\text{-}25$ & $02\text{-}15\text{-}34\ /\ 05\text{-}12\text{-}34$ \\
$18$ & $01\text{-}24\text{-}35\ /\ 02\text{-}13\text{-}45$ & $03\text{-}14\text{-}25$ & $04\text{-}15\text{-}23$ \\
$19$ & $01\text{-}24\text{-}35\ /\ 02\text{-}13\text{-}45$ & $03\text{-}14\text{-}25$ & $04\text{-}15\text{-}23\ /\ 05\text{-}12\text{-}34$ \\
$20$ & $01\text{-}24\text{-}35\ /\ 02\text{-}13\text{-}45$ & $03\text{-}14\text{-}25\ /\ 04\text{-}15\text{-}23$ & $05\text{-}12\text{-}34$ \\
$21$ & $01\text{-}23\text{-}45\ /\ 01\text{-}24\text{-}35\ /\ 02\text{-}13\text{-}45$ & $03\text{-}14\text{-}25$ & $05\text{-}12\text{-}34$ \\
$22$ & $01\text{-}25\text{-}34\ /\ 02\text{-}14\text{-}35\ /\ 03\text{-}12\text{-}45$ & $04\text{-}15\text{-}23$ & $05\text{-}13\text{-}24$ \\
$23$ & $01\text{-}24\text{-}35\ /\ 02\text{-}13\text{-}45\ /\ 02\text{-}14\text{-}35\ /\ 03\text{-}14\text{-}25$ & $04\text{-}15\text{-}23$ & $05\text{-}12\text{-}34$ \\
\bottomrule
\end{tabular}}
\end{center}

For $D=4$ colors, the $5$ orbit representatives (each class $P_i^k$
nonempty, classes pairwise edge-disjoint) are:

\begin{center}
\resizebox{\textwidth}{!}{%
\begin{tabular}{cllll}
\toprule
$k$ & $P_0^k$ & $P_1^k$ & $P_2^k$ & $P_3^k$ \\
\midrule
$0$ & $01\text{-}23\text{-}45$ & $02\text{-}14\text{-}35$ & $03\text{-}15\text{-}24$ & $04\text{-}13\text{-}25$ \\
$1$ & $01\text{-}23\text{-}45$ & $02\text{-}14\text{-}35$ & $03\text{-}15\text{-}24$ & $04\text{-}13\text{-}25\ /\ 05\text{-}12\text{-}34$ \\
$2$ & $01\text{-}23\text{-}45$ & $02\text{-}14\text{-}35$ & $03\text{-}15\text{-}24\ /\ 04\text{-}13\text{-}25$ & $05\text{-}12\text{-}34$ \\
$3$ & $01\text{-}23\text{-}45$ & $02\text{-}15\text{-}34\ /\ 03\text{-}14\text{-}25$ & $04\text{-}12\text{-}35$ & $05\text{-}13\text{-}24$ \\
$4$ & $01\text{-}24\text{-}35\ /\ 02\text{-}13\text{-}45$ & $03\text{-}14\text{-}25$ & $04\text{-}15\text{-}23$ & $05\text{-}12\text{-}34$ \\
\bottomrule
\end{tabular}}
\end{center}

Representative $0$ has $F=3$ free edges ($4{,}913=17^3$ assignments);
representatives $1$--$4$ have $F=0$ (their classes partition all
$15$ edges); total $4{,}917$ rows.

For $D=5$ colors the unique orbit representative is the ordered
$1$-factorization
\[
P_0=\{01\text{-}23\text{-}45\},\quad
P_1=\{02\text{-}14\text{-}35\},\quad
P_2=\{03\text{-}15\text{-}24\},\quad
P_3=\{04\text{-}13\text{-}25\},\quad
P_4=\{05\text{-}12\text{-}34\},
\]
with $F=0$: one assignment, eliminated by the conditions.

\section{Reproducibility}\label{app:repro}

All computations run on a standard Python 3.9 + numpy 2.0 + sympy 1.14
stack (z3 for the SAT cross-check), take a few seconds for the batch,
and contain no randomized components on the certification path.
Scripts: \path{computation/gen_search.py} (class enumeration and
batch), \path{computation/search_core.py} (exact decision procedure),
and the independent verification suite
\path{audit/verification_2026-09-03/} described in
Section~\ref{sec:verify}. The complete audit report is
\path{audit/AUDIT_REPORT_2026-09-03.md}.

\bibliographystyle{plain}
\begin{thebibliography}{9}

\bibitem{KGZ17} M. Krenn, X. Gu, A. Zeilinger,
\emph{Quantum experiments and graphs: Multiparty states as coherent
superpositions of perfect matchings}, Phys.\ Rev.\ Lett.\ 119, 240403
(2017).

\bibitem{KGS19} M. Krenn, X. Gu, D. Solt\'esz,
\emph{Questions on the structure of perfect matchings inspired by
quantum physics}, Proc.\ 2nd Croatian Combinatorial Days (2019),
57--70; arXiv:1902.06023.

\bibitem{bounty} M. Krenn, \emph{Inherited vertex coloring of graphs},
bounty page, \url{https://mariokrenn.wordpress.com/graph-theory-question/}
(mirrored verbatim at
\url{https://openproblemgarden.org/op/monochromatic_vertex_colorings_inherited_from_perfect_matchings})

\bibitem{GKM24} L. S. Chandran, R. Gajjala, C. Illickan,
\emph{Krenn--Gu conjecture for sparse graphs}, MFCS 2024, LIPIcs 306,
art.\ 41; arXiv:2407.00303.

\bibitem{CG24} L. S. Chandran, R. Gajjala,
\emph{Graph-theoretic insights on the constructability of complex
entangled states}, Quantum 8, 1396 (2024),
DOI 10.22331/q-2024-07-03-1396; arXiv:2304.06407.

\bibitem{CGI24} L. S. Chandran, R. Gajjala,
\emph{Edge-coloured graphs with only monochromatic perfect matchings
and their connection to quantum physics}, Electron.\ J.\ Combin.\
(2026); arXiv:2202.05562.

\bibitem{CLA22} A. Cervera-Lierta, M. Krenn, A. Aspuru-Guzik,
\emph{Design of quantum optical experiments with logic artificial
intelligence}, Quantum 6, 836 (2022); arXiv:2109.13273.

\bibitem{mathoverflow} MathOverflow 311325, \emph{A question on
inherited vertex colorings of perfect matchings} (M. Krenn, answers by
I. Bogdanov et al.).

\bibitem{alphaproof} DeepMind, \emph{Advancing mathematics research
with AI-driven formal proof search}, arXiv:2605.22763 (2026).

\bibitem{formalconjectures} google-deepmind/formal-conjectures,
\emph{MonochromaticQuantumGraph} Lean formalizations and open
statements, \url{https://github.com/google-deepmind/formal-conjectures}

\bibitem{notes} \emph{Krenn--Gu research workspace: proof notes and
certificate}, accompanying repository, 2026.

\bibitem{repo} \emph{Accompanying repository: scripts, logs, audit
report and certificate}, [URL to be added at submission].

\end{thebibliography}

\end{document}
